\documentclass[11pt,letterpaper]{article}
\usepackage[margin=1in]{geometry}
\usepackage[T1]{fontenc}
\usepackage{lmodern}
\usepackage{amsmath,amssymb,booktabs,tabularx,array}
\usepackage[hidelinks]{hyperref}

\setlength{\parindent}{0pt}
\setlength{\parskip}{0.4em}
\newcolumntype{Y}{>{\raggedright\arraybackslash}X}
\newcommand{\code}[1]{\texttt{#1}}

\begin{document}

\begin{center}
{\Large\bfseries Consolidating Partial Shipments into Truckload Tours\par}
\vspace{5pt}
MGSC 404 \enspace|\enspace Group 3, Project 1\\[3pt]
{\small Sofia Avetisian, Liana Jia, Abraham Dumani Lichtinger\\
Luca Novac, Vassili Prokopenko, Massimo Seifarth}
\end{center}

\section*{The planning objective}

We minimize the cost of serving each shipment on its required day by
dedicated truck, shared truck, or outside less-than-truckload (LTL)
service. We assume the outside carrier has enough capacity and can
meet the supplied windows at the workbook's rates. We develop the
model on 60 shipments over three days with eight
tractors at two depots, then extend it to 800 shipments over seven days
with 24 tractors at four depots. All demand is known before planning,
and every shipment is indivisible.

\section{Simplified workbook}

\subsection*{Parameters: sets, relationships and coefficients}

Python constructs the scheduled tour set and fixed costs before
Gurobi chooses which tours to run.

\begin{itemize}
\setlength{\itemsep}{2pt}
\item \(I,K,D\) are the shipments, tractors, and planning days.
\(R\) is the set of feasible candidate tours. Each tour fixes its
tractor, day, stop order, and times. \(R_{kd}\) contains the tours for
tractor \(k\) on day \(d\).
\item \(a_{ir}\) equals 1 if candidate \(r\) serves shipment \(i\)
and 0 otherwise. This is fixed once a tour is generated.
\item \(m_{ab},\tau_{ab}\) are miles and driving hours from location \(a\)
to \(b\). Service event \(j\) has handling time \(s_j\) and effective
window \([L_j,U_j]\). For tractor \(k\) on day \(d\),
\([H^-_{kd},H^+_{kd}]\) bounds departure from and return to its depot.
\item \(q_i^u,Q_k^u\) are shipment demand and trailer capacity for
resource \(u\), measured in pounds, feet of floor length, or pallets.
\item \(M_r\) is total tour mileage including depot legs, and \(N_r\)
counts pickups and deliveries. The equipment's minimum linehaul charge,
mileage rate, and extra-stop fee are \(\mu_r,\rho_r,f_r\).
\item \(b_i\) is the matched LTL base charge by freight class, weight,
and direct-mileage band. \(\phi_i\) is the fuel surcharge rate.
\end{itemize}

The resulting cost coefficients, in CAD, are
\[
c_r=\max\{\mu_r,\rho_r M_r\}+f_r\max\{0,N_r-2\},
\qquad \ell_i=b_i(1+\phi_i).
\]

\subsection*{Decision variables}

\begin{itemize}
\setlength{\itemsep}{2pt}
\item \textbf{Tour selection.} \(x_r=1\) if candidate tour \(r\)
is selected and 0 otherwise.
\item \textbf{Outsourcing.} \(z_i=1\) if shipment \(i\) is sent by LTL
and 0 otherwise.
\end{itemize}

\newpage
\subsection*{Mathematical model}

\begin{align}
\min\quad & \sum_{r\in R} c_r x_r+\sum_{i\in I}\ell_i z_i
   \tag{0}\label{eq:objective}\\
\text{subject to}\quad
& \sum_{r\in R}a_{ir}x_r+z_i=1 && \forall i\in I
   \tag{1}\label{eq:coverage}\\
& \sum_{r\in R_{kd}}x_r\le 1 && \forall k\in K,\ d\in D
   \tag{2}\label{eq:tractor}\\
& x_r,z_i\in\{0,1\} && \forall r\in R,\ i\in I.
   \tag{3}\label{eq:binary}
\end{align}

\subsection*{What each equation captures}

\begin{itemize}
\setlength{\itemsep}{1pt}
\item \eqref{eq:objective} \textbf{Total cost.} Add selected tour
charges and the LTL bill. Setting \(z_i=1\) for every shipment and
selecting no tours gives the all-LTL benchmark.
\item \eqref{eq:coverage} \textbf{Shipment coverage.} Each shipment
goes on exactly one selected tour or to LTL.
\item \eqref{eq:tractor} \textbf{Tractor availability.} Each tractor
can run at most one tour per day.
\item \eqref{eq:binary} \textbf{Domains.} Both choices are binary.
\end{itemize}

\subsection*{What makes a simplified tour feasible}

\textbf{Route and capacity.} A tour leaves and returns empty to its
home depot. Each shipment is picked up before delivery on the same
tractor and stated day. Side access allows delivery of any shipment
on board. A tour has at most six service events. After each event the
onboard load must satisfy
\[
0\le \sum_{i\text{ on board}}q_i^u\le Q_k^u
\qquad \text{for weight, floor length, and pallets}.
\]
The dry trailer holds 45,000 lb, 53 ft, and 26 pallets.

\textbf{Timing and driver hours.} Let \(A_j,T_j,E_j\) be arrival,
handling start, and departure at stop \(j\). The effective window
intersects the shipment window and facility hours. With \(E_0\) the
depot departure and \(A_{\rm home}\) the depot return,
\[
\begin{aligned}
A_j &= E_{j-1}+\tau_{j-1,j}, &
T_j &\ge \max\{A_j,L_j\},\\
T_j+s_j &\le U_j, &
E_j &\ge T_j+s_j .
\end{aligned}
\]
The tour fits depot hours and the shift when \(E_0\ge H^-_{kd}\)
and \(A_{\rm home}\le H^+_{kd}\). Driving, including both depot legs,
is capped at 8 hours. Duty \(A_{\rm home}-E_0\) includes driving,
handling, and waiting and is capped at 12 hours. Fixed shifts
provide at least 12 hours of rest.

\newpage
\section{Implementation}

\subsection*{Generate tours on the simplified workbook}

For each day, group tractors with the same depot, equipment, and shift.
Try groups of up to three shipments and all pickup-before-delivery
orders. Simulate each schedule and shift departure later when window
slack can absorb waiting. Keep the cheapest feasible schedule per group
and tractor, breaking ties by shorter duty. The retained schedules
supply \(R\), \(a_{ir}\), and \(c_r\) to Gurobi.

\subsection*{Extend the tour rules for the main workbook}

The same selection equations apply to the main workbook.

\textbf{Equipment and loading.} Chilled freight requires a reefer.
Shipments carried together must share temperature and compatibility
group. The dry trailer holds 45,000 lb, 53 ft, and 26 pallets, while
the reefer holds 42,000 lb, 48 ft, and 24 pallets. Under last-in,
first-out unloading, the most recently picked-up shipment still on
board must be delivered first. The limit rises to eight service
events, or four shipments.

\textbf{Time windows and driver limits.} The same time-window formulas
use each location's opening hours. Driving is capped at 11 hours
and duty at 14 hours, with at least 10 hours of rest between days.
Before a leg would take the driver past eight driving hours since
departure or the last break, the driver must take a 30-minute break
at a service stop. The break fits around handling and counts toward
duty. Driving back to the depot counts toward both driving limits.

\textbf{Candidate search.} Python extends a group only when every
immediate subset yielded a feasible tour for that tractor profile.
It screens stop orders for capacity, loading order, and early window
violations. If the initial schedule fails on a route that needs a
break, it tries alternative break placements and departures on the
15-minute grid. The run allows up to four shipments per tour with
no additional cap on the number of groups tried. The solver gap
applies to the candidates generated by this search.

\subsection*{Solve, export and check}

From the project folder, run \code{python Gurobi\_Simplified\_v4.py}
or \code{python Gurobi\_Main\_V1.py}. Each script reads its workbook,
exports the selected tours, stop times, loads, and LTL assignments,
then calls the checker. Both runs used Gurobi 13.0.3, seed 0, and a
target gap of 0\%. Time limits were 600 and 1,800 seconds.

The checker reads the workbook and exported plan independently of
Gurobi. It rebuilds the stop sequence, travel, waiting, loads, return
to depot, and costs, then checks coverage and driver limits. Both
exported plans passed.
Running \code{python checker-simplified.py --self-test} verified
three rejections on temporary copies. The checker flagged S020 at
45,001 lb for exceeding capacity, rejected its late pickup for missing
the service window, and flagged the second Day 1 tour assigned to T01.

\newpage
\section{Results and managerial implications}

\textbf{Cost.} The simplified plan costs \textbf{CAD 36,246.15},
saving CAD 57,052.20 (61.15\%) against all-LTL. The main plan costs
\textbf{CAD 1,137,319.80}, saving CAD 420,216.05 (26.98\%).
The main bill comprises CAD 178,739.60 for fleet tours and
CAD 958,580.20 for LTL. Appendix~\ref{app:results} gives the
operating results and savings-heuristic comparison.

\textbf{Worked example from the simplified plan.} On Day 1, T01 runs
the appendix tour for S020 and S017. For this tour \(r^\star\),
\(a_{\mathrm{S020},r^\star}=a_{\mathrm{S017},r^\star}=1\). Selecting
\(x_{r^\star}=1\) makes both LTL decisions zero in \eqref{eq:coverage}
and uses T01's Day 1 slot in \eqref{eq:tractor}. The tour travels
303 miles and has four service events. The dry trailer's cost is
\[
c_{r^\star}=\max\{900,2.20(303)\}+50(4-2)=1{,}000.
\]
The two LTL charges total CAD 2,761.15, so the tour saves CAD 1,761.15.
Peak load is 26,000 lb, 36 ft, and 14 pallets. Across the simplified
plan, 24 tours serve 52 shipments and eight go by LTL. S007 has no
generated fleet option. S008 has positive-saving candidates competing
for tractor-days already assigned to other tours.

\textbf{Main fleet deployment.} Shared tours handle 273 shipments,
dedicated tours handle 38, and LTL handles 489. All 168 tractor-days
are used, including all 28 reefer-days. The fleet carries 17 of 78
chilled shipments. Refrigeration limits tractor choice, and
compatibility and loading order restrict which shipments can share
a trailer and in what order they can be delivered. Varied opening
hours narrow the timing options, and 111 selected tours schedule a
driver break.

\textbf{Trailer use and service.} In the main plan, weight utilization
is 16.91\% and floor utilization is 22.22\%, including empty travel.
Empty legs account for 37,298 of 68,388 fleet miles (54.5\%).
Elapsed time from pickup completion to delivery completion averages
about 17 minutes above the direct-trip benchmark, with a four-hour
maximum. All fleet shipments are picked up and delivered on their
stated day.

\textbf{Recommendation.} Use shared tours where the avoided LTL
charges cover the extra stops and detours. Keep dedicated tours
available when their savings justify a tractor-day. Use LTL when
the fleet has no suitable candidate or another assignment saves
more. When comparing tours, consider depot approach and return
mileage alongside weight and floor utilization.

\textbf{Scope.} Gurobi reported \code{OPTIMAL} with a zero gap to
reported precision for the 1,592 simplified and 12,533 main
candidates. The independent checker verified the selected plans
and charges. The savings reflect the supplied rates and demand.

\clearpage
\appendix
\section{Results for the supplied workbooks}\label{app:results}

\subsection*{Operating results and benchmark}

\begin{center}
\small
\begin{tabularx}{\linewidth}{@{}Yrr@{}}
\toprule
Measure & Simplified & Main\\
\midrule
All-LTL cost (CAD) & 93,298.35 & 1,557,535.85\\
Savings heuristic cost (CAD) & 41,401.80 & 1,227,702.10\\
Gurobi cost below heuristic (CAD) & 5,155.65 & 90,382.30\\
Shipments on shared / dedicated tours / LTL & 52 / 0 / 8 & 273 / 38 / 489\\
Tractor-days used / available & 24 / 24 & 168 / 168\\
Weight / floor utilization (\%) & 27.40 / 30.22 & 16.91 / 22.22\\
Tour miles, including empty legs & 7,698 & 68,388\\
Empty miles (depot legs only) & 3,242 & 34,953\\
Empty miles (all empty legs) & 3,242 & 37,298\\
Added elapsed hours, mean / maximum & 1.83 / 6.50 & 0.29 / 4.00\\
Candidate tours & 1,592 & 12,533\\
Status / gap over candidates & Optimal / 0.00\% & Optimal / 0.00\%\\
Independent feasibility check & Pass & Pass\\
Generation / Gurobi time (s) & 2.464 / 0.040 & 1.186 / 0.200\\
\bottomrule
\end{tabularx}
\end{center}

For the savings heuristic, use the model's feasible schedules for
two-shipment tours. Rank pairs by positive savings versus LTL, rounded
to cents, with ties broken by day, shipment IDs, then tractor ID. Take
each pair in order if both shipments and its tractor-day are free. Send
the rest by LTL. Both model scripts report the benchmark.

Utilization is mile-weighted and includes empty legs. Added elapsed
time runs from pickup completion to delivery completion, less direct
driving and the same delivery handling. The direct benchmark excludes
waiting for the delivery window.

\subsection*{Example simplified tour, Day 1, tractor T01}

\begin{center}
\small
\begin{tabularx}{\linewidth}{@{}llYr@{}}
\toprule
Local time & Location & Event & Onboard lb / ft / pallets\\
\midrule
06:30 & D01 & Depart & 0 / 0 / 0\\
09:00--09:30 & F09 & Pick up S020 & 15,000 / 18 / 6\\
10:00--10:30 & F05 & Pick up S017 & 26,000 / 36 / 14\\
12:00--12:30 & F08 & Deliver S017 & 15,000 / 18 / 6\\
13:30--14:00 & F01 & Deliver S020 & 0 / 0 / 0\\
15:00 & D01 & Return & 0 / 0 / 0\\
\bottomrule
\end{tabularx}
\end{center}

\end{document}
